\chapter{Linear Algebra and Calculus}\label{ch:iv:linear-algebra-and-calculus}

A risk system rejects the morning's correlation matrix: one eigenvalue is $-0.8$. The quant on call is asked, before
touching any code, why three pairwise correlations that each look reasonable can be impossible together, and what
the nearest valid matrix is. Linear algebra and calculus questions in interviews are rarely about technique for its
own sake; they are about the few facts that recur in finance (a covariance matrix must be positive semidefinite, a
regression is a projection, a Gaussian integral has a closed form, a constrained optimum has a multiplier) and about
using them quickly.

\section{Eigenvalues, definiteness and correlation matrices}

A covariance or correlation matrix is symmetric positive semidefinite: $w^\top \Sigma w$ is the variance of the
portfolio $w$ and cannot be negative (One Quant Book~4, chapter~22). Equivalently all eigenvalues are non-negative,
or all principal minors are. Pairwise estimation, missing data and hand-edited stress scenarios break this
property, and interviews ask how to see it and repair it.

\begin{proposition}[The third correlation]\label{prop:iv:linear-algebra-and-calculus:third}
Given $\rho_{12}$ and $\rho_{13}$, the $3 \times 3$ correlation matrix is positive semidefinite exactly when
\[
\rho_{12}\rho_{13} - \sqrt{(1 - \rho_{12}^2)(1 - \rho_{13}^2)} \;\le\; \rho_{23} \;\le\;
\rho_{12}\rho_{13} + \sqrt{(1 - \rho_{12}^2)(1 - \rho_{13}^2)}.
\]
\end{proposition}

\begin{proof}
The determinant $1 - \rho_{12}^2 - \rho_{13}^2 - \rho_{23}^2 + 2\rho_{12}\rho_{13}\rho_{23}$ is a concave quadratic in
$\rho_{23}$ whose roots are the two bounds; between them it is non-negative, and the $2 \times 2$ minors are
non-negative for correlations in $[-1, 1]$.
\end{proof}

\begin{omfigure}
\begin{tikzpicture}
\begin{axis}[width=9.5cm,height=5.6cm,xlabel={\small $\rho_{12}$},ylabel={\small feasible $\rho_{23}$},
  xmin=-1,xmax=1,ymin=-1,ymax=1,xtick={-1,-0.5,0,0.5,1},
  tick label style={font=\footnotesize},grid=major,grid style={gray!20},table/col sep=comma,
  legend cell align=left,legend style={font=\scriptsize,at={(0.5,-0.3)},anchor=north,
  /tikz/every even column/.append style={column sep=8pt}},legend columns=2]
\addplot[omDef,thick] table[x=r12,y=lo05]{figdata/interviews/15-linear-algebra-and-calculus/corr.csv};
\addplot[omThm,thick,dashed] table[x=r12,y=lo09]{figdata/interviews/15-linear-algebra-and-calculus/corr.csv};
\addplot[omDef,thick,forget plot] table[x=r12,y=hi05]{figdata/interviews/15-linear-algebra-and-calculus/corr.csv};
\addplot[omThm,thick,dashed,forget plot] table[x=r12,y=hi09]{figdata/interviews/15-linear-algebra-and-calculus/corr.csv};
\legend{{$\rho_{13} = 0.5$},{$\rho_{13} = 0.9$}}
\end{axis}
\end{tikzpicture}

\omcaption{The feasible values of $(\rho_{12}, \rho_{23})$ fill the inside of each ellipse
(\cref{prop:iv:linear-algebra-and-calculus:third}). When $\rho_{13} = 0.9$ the ellipse is thin: for a given $\rho_{12}$
the range of $\rho_{23}$ is narrow, and two assets each highly correlated with a third must be correlated with each
other. Data: \texttt{fig\_iv\_corr.py}.}\label{fig:iv:linear-algebra-and-calculus:third}
\end{omfigure}

\begin{example}[Equicorrelation]\label{ex:iv:linear-algebra-and-calculus:equi}
An $n \times n$ matrix with ones on the diagonal and $\rho$ elsewhere is $(1 - \rho)I + \rho\mathbf 1\mathbf 1^\top$. Its
eigenvalues are $1 + (n - 1)\rho$, with eigenvector $\mathbf 1$, and $1 - \rho$ with multiplicity $n - 1$, so it is a
correlation matrix exactly when $-1/(n-1) \le \rho \le 1$. Ten assets cannot all be pairwise correlated at $-0.2$.
\end{example}

\begin{method}[Repairing a correlation matrix]\label{met:iv:linear-algebra-and-calculus:repair}
\begin{enumerate}
\item Check: compute the smallest eigenvalue; a Cholesky factorisation that fails is the fast test.
\item Repair by clipping: set negative eigenvalues to zero (or a small floor), rebuild, and rescale to a unit
diagonal (eigenvalue clipping, One Quant Book~4, chapter~22).
\item Repair by the nearest correlation matrix in the Frobenius norm (Higham, 2002): alternate projections onto the
positive semidefinite cone and onto unit-diagonal matrices.
\item Say which entries moved most, since they are the ones the data did not support.
\end{enumerate}
\end{method}

\section{Projections and least squares}

Ordinary least squares projects the vector of observations $y$ onto the column space of the regressors $X$. The hat
matrix $H = X(X^\top X)^{-1}X^\top$ is symmetric and idempotent, its eigenvalues are 0 and 1, and its trace is the number
of regressors $p$; its diagonal entries are the leverages, which average $p/n$. The residual $y - Hy$ is orthogonal to
every column of $X$, and $R^2$ is the squared cosine of the angle between the centred $y$ and its projection.

\begin{example}[A rank-one update]\label{ex:iv:linear-algebra-and-calculus:rankone}
$I + vv^\top$ has eigenvalue $1 + \|v\|^2$ in the direction of $v$ and 1 on the orthogonal complement, and its
inverse is $I - vv^\top/(1 + \|v\|^2)$ (Sherman--Morrison). A one-factor covariance matrix $\sigma^2 I + \beta\beta^\top$
is this form, and its inverse is what minimum-variance portfolios need.
\end{example}

\section{Integrals, series and expansions}

A handful of results covers most calculus questions: $\int e^{-x^2/2}\,dx = \sqrt{2\pi}$; for a standard normal $Z$,
$\E[e^{aZ}] = e^{a^2/2}$, $\E[Z^4] = 3$ and $\E[Z^+] = 1/\sqrt{2\pi}$; $\sum_{k \ge 1} k x^k = x/(1-x)^2$ for $|x| < 1$;
and Taylor's theorem with a remainder, which turns ``approximately'' into a bound.

\begin{method}[Integrals in an interview]\label{met:iv:linear-algebra-and-calculus:integrals}
\begin{enumerate}
\item Look for a density: an integrand that is a density times something is an expectation.
\item Complete the square in an exponent; it turns a Gaussian integral with a linear term into a shifted one.
\item Use symmetry to kill odd moments.
\item Differentiate under the integral sign, or with respect to a parameter of a known series.
\end{enumerate}
\end{method}

\section{Optimisation with a constraint}

A constrained optimum sets the gradient of the objective parallel to the gradient of the constraint. The minimum-variance
portfolio fully invested ($\mathbf 1^\top w = 1$) is $w = \Sigma^{-1}\mathbf 1/(\mathbf 1^\top\Sigma^{-1}\mathbf 1)$; the
portfolio of highest expected return at a given volatility is proportional to $\Sigma^{-1}\mu$, scaled to the
volatility (One Quant Book~7, chapter~25, on mean--variance construction).

\begin{example}[Two assets]\label{ex:iv:linear-algebra-and-calculus:twoasset}
With volatilities $\sigma_1, \sigma_2$ and correlation $\rho$, the minimum-variance weight on the first asset is
$(\sigma_2^2 - \rho\sigma_1\sigma_2)/(\sigma_1^2 + \sigma_2^2 - 2\rho\sigma_1\sigma_2)$. With 20\%, 10\% and 0.3 it is
about 0.105, for a portfolio volatility of about 9.8\%, below that of either asset.
\end{example}

\section{Worked answers}

\begin{example}[The volatility of a pair trade]\label{ex:iv:linear-algebra-and-calculus:pair}
\emph{``Long one unit of a stock with 20\% volatility, short one unit of a peer with 25\%, correlation 0.9. What is the
volatility of the pair?''} Write the quadratic form $w^\top\Sigma w$ with $w = (1, -1)$: $0.2^2 + 0.25^2 - 2 \times 0.9
\times 0.2 \times 0.25 = 0.04 + 0.0625 - 0.09 = 0.0125$, a volatility of $\sqrt{0.0125} \approx 11.2\%$. \emph{Checks:} at
correlation 1 the pair's volatility would be the difference, 5\%; at 0 it would be $\sqrt{0.1025} \approx 32\%$; 11\% lies
between and near the high-correlation end. The answer also shows why pairs are sized by volatility: with the short
scaled to $0.2/0.25 = 0.8$ units the variance falls to $0.04 + 0.04 - 2 \times 0.9 \times 0.04 = 0.008$, about 8.9\%.
\end{example}

\begin{example}[Integration by parts against a normal]\label{ex:iv:linear-algebra-and-calculus:stein}
\emph{``For a standard normal $Z$, compute $\E[Z\sin Z]$.''} Since $\varphi'(z) = -z\varphi(z)$, integrating by parts gives
$\E[Zf(Z)] = \E[f'(Z)]$ for smooth $f$ of moderate growth. So $\E[Z\sin Z] = \E[\cos Z]$, the real part of the
characteristic function at 1: $e^{-1/2} \approx 0.607$. The identity is worth knowing by name (Stein's lemma), because it
also gives $\E[Z^4] = 3\E[Z^2] = 3$ in one line and appears in the Greeks of \cref{ch:iv:options-and-derivatives}.
\end{example}

\begin{example}[Newton's method on a transcendental equation]\label{ex:iv:linear-algebra-and-calculus:newton}
\emph{``Solve $xe^x = 1$ to four decimals, by hand.''} With $g(x) = xe^x - 1$ and $g'(x) = (1 + x)e^x$, start at $x_0 =
0.5$, where $g = 0.5 \times 1.6487 - 1 = -0.1756$: the step is $0.1756/(1.5 \times 1.6487) \approx 0.071$, so $x_1
\approx 0.5710$. One more step gives $x_2 \approx 0.56716$, and the next changes only the sixth decimal: $0.5671$. Newton
roughly doubles the number of correct digits at each step, which is why the same method solves for an implied volatility
or a yield in two or three iterations from a sensible start.
\end{example}

\section{Question bank}

\begin{interviewq}[$\star$ \iqroles{researcher, bank} \iqfirm{any}]\label{iq:iv:linear-algebra-and-calculus:1}
What are the eigenvalues and eigenvectors of $\begin{pmatrix} 2 & 1 \\ 1 & 2 \end{pmatrix}$?
\end{interviewq}

\begin{interviewq}[$\star$ \iqroles{researcher, risk} \iqfirm{bank}]\label{iq:iv:linear-algebra-and-calculus:2}
Is the matrix with unit diagonal, $\rho_{12} = \rho_{13} = 0.9$ and $\rho_{23} = -0.9$ a valid correlation matrix? How do
you see it quickly?
\end{interviewq}

\begin{interviewq}[$\star$ \iqroles{researcher, bank} \iqfirm{bank}]\label{iq:iv:linear-algebra-and-calculus:3}
For a standard normal $Z$, compute $\E[e^Z]$ and $\E[Z^4]$.
\end{interviewq}

\begin{interviewq}[$\star$ \iqroles{researcher, trader} \iqfirm{any}]\label{iq:iv:linear-algebra-and-calculus:4}
Compute $\sum_{k \ge 1} k/2^k$.
\end{interviewq}

\begin{interviewq}[$\star\star$ \iqroles{risk, researcher} \iqfirm{bank}]\label{iq:iv:linear-algebra-and-calculus:5}
Two assets have correlations 0.8 and 0.6 with a third. What values can their correlation with each other take?
\end{interviewq}

\begin{interviewq}[$\star\star$ \iqroles{researcher} \iqfirm{systematic fund}]\label{iq:iv:linear-algebra-and-calculus:6}
Ten assets are all pairwise correlated at $\rho$. What are the eigenvalues of the correlation matrix, and what is the
smallest possible $\rho$?
\end{interviewq}

\begin{interviewq}[$\star\star$ \iqroles{researcher, mle} \iqfirm{systematic fund}]\label{iq:iv:linear-algebra-and-calculus:7}
In a regression of 250 observations on an intercept and four regressors, what are the trace and the eigenvalues of the
hat matrix? What is the average leverage, and why does it matter?
\end{interviewq}

\begin{interviewq}[$\star\star$ \iqroles{risk, trader} \iqfirm{asset manager}]\label{iq:iv:linear-algebra-and-calculus:8}
Two assets have volatilities 20\% and 10\% and correlation 0.3. What fully invested portfolio has the lowest variance,
and what is its volatility?
\end{interviewq}

\begin{interviewq}[$\star\star$ \iqroles{bank, researcher} \iqfirm{bank}]\label{iq:iv:linear-algebra-and-calculus:9}
For a standard normal $Z$, compute $\E[\max(Z, 0)]$. Where does this number appear in option pricing?
\end{interviewq}

\begin{interviewq}[$\star\star\star$ \iqroles{risk, developer} \iqfirm{bank}]\label{iq:iv:linear-algebra-and-calculus:10}
Find the valid correlation matrix nearest, in the Frobenius norm, to the invalid matrix of
\cref{iq:iv:linear-algebra-and-calculus:2}. How far does it move, and which algorithm scales to a thousand assets?
\end{interviewq}

\begin{interviewq}[$\star\star\star$ \iqroles{researcher, bank} \iqfirm{systematic fund}]\label{iq:iv:linear-algebra-and-calculus:11}
With $v = (1, 2, 2)^\top$, what are the eigenvalues of $I + vv^\top$, and what is its inverse? Why does this matter for a
one-factor covariance model?
\end{interviewq}

\begin{interviewq}[$\star\star\star$ \iqroles{researcher, trader} \iqfirm{multi-manager fund}]\label{iq:iv:linear-algebra-and-calculus:12}
Two uncorrelated assets have expected returns 5\% and 3\% and volatilities 20\% and 10\%. Which portfolio has the highest
expected return at a volatility of 10\%, with no budget constraint? What are its weights, its expected return and its
Sharpe ratio?
\end{interviewq}

\begin{interviewq}[$\star\star\star$ \iqroles{bank, researcher} \iqfirm{any}]\label{iq:iv:linear-algebra-and-calculus:13}
Bound the error of $1 + x + x^2/2$ as an approximation of $e^x$ at $x = 0.1$, and compare with the actual error.
\end{interviewq}

\begin{omsources}
\item N.~J.~Higham, ``Computing the nearest correlation matrix---a problem from finance'', \emph{IMA Journal of
Numerical Analysis} 22(3), 2002, 329--343.
\item One Quant Book~4, chapters 16, 22 and~25 (least squares, covariance estimation, numerical linear algebra); One
Quant Book~7, chapter~25 (portfolio construction).
\end{omsources}
